\documentclass[../../main.tex]{subfiles}
\graphicspath{{\subfix{../../image/}}}

\begin{document}

\section{不变性，Noether 定理}\label{section:ch08-6-不变性Noether}

接下来研究在适当的定义域与函数变分下不变的变分被积函数，并证明相应 Euler--Lagrange 方程的解自动也求解某些散度结构守恒律。

\subsection{不变变分问题}\label{subsection:ch08-6-1-不变变分问题}

再次把注意力转向泛函
\begin{align}
I[w]\triangleq\int_U L(Dw,w,x)\,\mathrm{d}x,\label{align:ch08-6-eq01}
\end{align}
其中 $U\subset\mathbb{R}^n$，$w:U\to\mathbb{R}$。如常写 $L=L(p,z,x)$。

\paragraph{记号。}
(i) 设 $\mathbf{x}:\mathbb{R}^n\times\mathbb{R}\to\mathbb{R}^n$，$\mathbf{x}=\mathbf{x}(x,\tau)$，为光滑向量场族，满足 $\mathbf{x}(x,0)=x$ 对所有 $x\in\mathbb{R}^n$。则对小 $|\tau|$，映射 $x\mapsto\mathbf{x}(x,\tau)$ 是光滑微分同胚。称映射 $x\mapsto\mathbf{x}(x,\tau)$ 为\textbf{定义域变分}。还定义
\begin{align}
\mathbf{v}(x)\triangleq\mathbf{x}_\tau(x,0)\label{align:ch08-6-eq02}
\end{align}
及
\begin{align}
U(\tau)\triangleq\mathbf{x}(U,\tau).\label{align:ch08-6-eq03}
\end{align}
(ii) 接着，给定光滑 $u:\mathbb{R}^n\to\mathbb{R}$，考虑光滑函数变分族 $w:\mathbb{R}^n\times\mathbb{R}\to\mathbb{R}$，$w=w(x,\tau)$，满足
\begin{align}
w(x,0)=u(x).\label{align:ch08-6-eq04}
\end{align}
写
\begin{align}
m(x)\triangleq w_\tau(x,0).\label{align:ch08-6-eq05}
\end{align}
由于稍后将清楚的原因，称 $m$ 为\textbf{乘子}。

给定形如\eqref{align:ch08-6-eq01} 的泛函 $I[\cdot]$，我们问：能否找到与 Lagrangian $L$ 相容的定义域变分与函数变分，即 $I[\cdot]$ 在这些变分下不变？

\begin{definition}\label{definition:ch08-6-不变}
称泛函 $I[\cdot]$ 在定义域变分 $\mathbf{x}$ 与函数变分 $w$ 下\textbf{不变}，只要
\begin{align}
\int_U L(Dw(x,\tau),w(x,\tau),x)\,\mathrm{d}x=\int_{U(\tau)}L(Du,u,x)\,\mathrm{d}x\label{align:ch08-6-eq06}
\end{align}
对所有小 $|\tau|$ 及所有开集 $U\subset\mathbb{R}^n$ 成立。
\end{definition}

此处写 $Dw=D_xw$。此定义背后的想法是：给定定义域变分 $\mathbf{x}$ 与函数 $u$，我们将把 $w$ 寻找为涉及 $u(\mathbf{x}(x,\tau))$ 的某个表达式。然后通过在左端积分中作变量替换来验证\eqref{align:ch08-6-eq06}，积分之后将在区域 $U(\tau)$ 上进行。

\subsection{Noether 定理}\label{subsection:ch08-6-2-Noether定理}

现证明：泛函的不变性蕴涵相应 Euler--Lagrange 方程可化为散度形式。

\begin{theorem}[Noether 定理]\label{theorem:ch08-6-1-Noether}
设泛函 $I[\cdot]$ 在对应于光滑函数 $u$ 的定义域变分 $\mathbf{x}$ 与函数变分 $w$ 下不变。
(i) 则
\begin{align}
\sum_{i=1}^n\big(mL_{p_i}(Du,u,x)-L(Du,u,x)v_i\big)_{x_i}=m\left(\sum_{i=1}^n(L_{p_i}(Du,u,x))_{x_i}-L_z(Du,u,x)\right),\label{align:ch08-6-eq07}
\end{align}
其中 $\mathbf{v}=(v_1,\cdots,v_n)$ 由\eqref{align:ch08-6-eq02} 定义，乘子 $m$ 由\eqref{align:ch08-6-eq05} 定义。
(ii) 特别地，若 $u$ 是 $I[\cdot]$ 的临界点从而求解 Euler--Lagrange 方程 $-\operatorname{div}(D_pL)+L_z=0$，则有散度恒等式
\begin{align}
\sum_{i=1}^n\big(mL_{p_i}(Du,u,x)-L(Du,u,x)v_i\big)_{x_i}=0.\label{align:ch08-6-eq08}
\end{align}
因此乘以 $m$ 把 Euler--Lagrange PDE 化为散度形式。
\begin{proof}
对不变性恒等式\eqref{align:ch08-6-eq06} 关于 $\tau$ 求导，再令 $\tau=0$，得恒等式
\begin{align*}
\int_U D_pL\cdot Dm+L_zm\,\mathrm{d}x=\int_{\partial U}L\mathbf{v}\cdot\nu\,\mathrm{d}S.
\end{align*}
右端项的出现是由于 \S C.4 中关于在运动区域上积分求导的公式。现分部积分并用 Gauss--Green 定理得
\begin{align*}
\int_U(-\operatorname{div}D_pL+L_z)m\,\mathrm{d}x&=\int_{\partial U}(-mD_pL+L\mathbf{v})\cdot\nu\,\mathrm{d}S\\
&=\int_U\operatorname{div}(-mD_pL+L\mathbf{v})\,\mathrm{d}x.
\end{align*}
此恒等式对所有区域 $U$ 成立，故恒等式\eqref{align:ch08-6-eq07} 得证。
\end{proof}
\end{theorem}

如前所述，有时可先猜出一个定义域变分 $\mathbf{x}$，然后寻找相应的函数变分 $w$ 作为涉及 $u(\mathbf{x}(x,\tau))$ 的某个公式。于是可用 $u$ 及其偏导数计算乘子 $m$。下列例子在越来越复杂的情形中说明这一程序。

\begin{example}[与 $x$ 无关的 Lagrangian]\label{example:ch08-6-1-平移不变}
若 $L=L(p,z)$ 不依赖于自变量 $x$，则积分\eqref{align:ch08-6-eq01} 在空间平移下不变。具体地，选 $k\in\{1,\cdots,n\}$，定义
\begin{align*}
\mathbf{x}(x,\tau)\triangleq x+\tau e_k,\qquad w(x,\tau)\triangleq u(x+\tau e_k).
\end{align*}
则
\begin{align*}
\mathbf{v}=e_k,\qquad m=u_{x_k}.
\end{align*}
因此若 $u$ 是临界点，\eqref{align:ch08-6-eq08} 成为恒等式
\begin{align}
\sum_{i=1}^n\big(L_{p_i}u_{x_k}-L\delta_{ik}\big)_{x_i}=0\qquad (k=1,\cdots,n).\label{align:ch08-6-eq09}
\end{align}
简单练习可验证这些公式直接由 Euler--Lagrange 方程推出。关键在于 Noether 定理为寻找此类恒等式提供了系统程序。
\end{example}

\begin{example}[伸缩不变性]\label{example:ch08-6-2-伸缩不变}
泛函
\begin{align*}
I[w]=\int_U|Dw|^p\,\mathrm{d}x
\end{align*}
（其光滑极小元 $u$ 求解 $p$-Laplace 方程 $\operatorname{div}(|Du|^{p-2}Du)=0$）在伸缩变换
\begin{align*}
x\mapsto\lambda x,\qquad u\mapsto\lambda^{\frac{n-p}{p}}u(\lambda x)
\end{align*}
下不变。为与前面记号一致，令 $\lambda=e^\tau$，定义
\begin{align*}
\mathbf{x}(x,\tau)\triangleq e^\tau x,\qquad w(x,\tau)\triangleq e^{\tau\frac{n-p}{p}}u(e^\tau x).
\end{align*}
则
\begin{align*}
\mathbf{v}=x,\qquad m=Du\cdot x+\frac{n-p}{p}u.
\end{align*}
$p$-Laplace 方程解的相应散度恒等式\eqref{align:ch08-6-eq08} 为
\begin{align}
\sum_{i=1}^n\left(\left(Du\cdot x+\frac{n-p}{p}u\right)p|Du|^{p-2}u_{x_i}-|Du|^p x_i\right)_{x_i}=0.\label{align:ch08-6-eq10}
\end{align}
直接验证此恒等式同样容易。
\end{example}

\paragraph{应用：单调性公式。}设 $u$ 是某区域 $U$ 内 $p$-Laplace PDE 的光滑解，且球 $B(0,r)\subset U$。若把散度恒等式\eqref{align:ch08-6-eq10} 在 $B(0,r)$ 上积分并用 Gauss--Green 定理化简，发现
\begin{align*}
(n-p)\int_{B(0,r)}|Du|^p\,\mathrm{d}x=r\int_{\partial B(0,r)}|Du|^p-p|Du|^{p-2}u_r^2\,\mathrm{d}S,
\end{align*}
其中 $u_r\triangleq\frac{Du\cdot x}{|x|}$。因此
\begin{align*}
\frac{\mathrm{d}}{\mathrm{d}r}\left(\frac1{r^{n-p}}\int_{B(0,r)}|Du|^p\,\mathrm{d}x\right)&=\frac{p-n}{r^{n-p+1}}\int_{B(0,r)}|Du|^p\,\mathrm{d}x+\frac1{r^{n-p}}\int_{\partial B(0,r)}|Du|^p\,\mathrm{d}S\\
&=\frac{p}{r^{n-p}}\int_{\partial B(0,r)}|Du|^{p-2}u_r^2\,\mathrm{d}S\geqslant 0.
\end{align*}
这是一个单调性公式，意指
\begin{align*}
r\mapsto\frac1{r^{n-p}}\int_{B(0,r)}|Du|^p\,\mathrm{d}x\text{ 非减}.
\end{align*}
我们在习题 20 中讨论 $p=2$ 时一个有趣的推论。
在 \S9.4.2 中将看到由例 2 提示的乘子 $Du\cdot x$ 在研究某些半线性椭圆方程时的用处。

\paragraph{时间相关问题。}若自变量之一被识别为时间，则可把\eqref{align:ch08-6-eq08} 解释为由变分积分不变性导致的守恒律。

\begin{example}[非线性波方程的能量守恒]\label{example:ch08-6-3-能量守恒}
考虑积分表达式
\begin{align}
I[w]=\int_0^T\int_{\mathbb{R}^n}\frac12w_t^2-\left(\frac12|Dw|^2+F(w)\right)\,\mathrm{d}x\,\mathrm{d}t,\label{align:ch08-6-eq11}
\end{align}
定义于（例如）具有紧支集的函数 $w=w(x,t)$。如常写 $Dw=D_xw$。与 \S3.3.1 的 ODE 例类比，可把 Lagrangian 解释为动能减势能。相应 Euler--Lagrange 方程是半线性波方程
\begin{align}
u_{tt}-\Delta u+f(u)=0,\label{align:ch08-6-eq12}
\end{align}
其中 $f\triangleq F'$。

\eqref{align:ch08-6-eq11} 的被积函数不依赖于时间变量 $t$，故在该变量平移下不变。Noether 定理蕴涵此不变性强制一个守恒律，在此情形即能量守恒。更精确地，定义
\begin{align*}
\mathbf{x}(x,t,\tau)\triangleq(x,t+\tau),\qquad w(x,t,\tau)\triangleq u(x,t+\tau),
\end{align*}
于是
\begin{align*}
\mathbf{v}=e_{n+1},\qquad m=u_t.
\end{align*}
则\eqref{align:ch08-6-eq08} 读为
\begin{align*}
\sum_{i=1}^n(-u_{x_i}u_t)_{x_i}+\left(u_t^2-\frac12(u_t^2-|Du|^2)+F(u)\right)_t=0.
\end{align*}
可改写为
\begin{align}
e_t-\operatorname{div}(u_tDu)=0,\label{align:ch08-6-eq13}
\end{align}
其中能量密度
\begin{align*}
e\triangleq\frac12(u_t^2+|Du|^2)+F(u).
\end{align*}
\eqref{align:ch08-6-eq13} 中的散度算子仅作用于 $x$ 变量。若 $u$ 在每个时刻在空间中具有紧支集，则总能量在时间中守恒：
\begin{align*}
\frac{\mathrm{d}}{\mathrm{d}t}\int_{\mathbb{R}^n}\frac12(u_t^2+|Du|^2)+F(u)\,\mathrm{d}x=0.
\end{align*}
在第 12 章将系统研究形如\eqref{align:ch08-6-eq12} 的半线性波方程的可解性，并将发现这一守恒律极为有用。
\end{example}

\begin{example}[波方程的伸缩不变性]\label{example:ch08-6-4-波方程伸缩}
线性波方程
\begin{align*}
\Box u=u_{tt}-\Delta u=0
\end{align*}
对应作用泛函
\begin{align}
I[w]=\frac12\int_0^T\int_{\mathbb{R}^n}w_t^2-|Dw|^2\,\mathrm{d}x\,\mathrm{d}t.\label{align:ch08-6-eq14}
\end{align}
与上面例 2 类似，泛函\eqref{align:ch08-6-eq14} 在伸缩变换
\begin{align*}
(x,t)\mapsto(\lambda x,\lambda t),\qquad u\mapsto\lambda^{\frac{n-1}{2}}u(\lambda x,\lambda t)
\end{align*}
下不变。如前令 $\lambda=e^\tau$，定义
\begin{align*}
\mathbf{x}(x,t,\tau)\triangleq(e^\tau x,e^\tau t),\qquad w(x,t,\tau)\triangleq e^{\tau\frac{n-1}{2}}u(e^\tau x,e^\tau t).
\end{align*}
则
\begin{align*}
\mathbf{v}=(x,t),\qquad m=tu_t+x\cdot Du+\frac{n-1}{2}u.
\end{align*}
\eqref{align:ch08-6-eq08} 给出的守恒律断言
\begin{align*}
p_t-\operatorname{div}\mathbf{q}=0,
\end{align*}
其中
\begin{align*}
p\triangleq\frac t2(u_t^2+|Du|^2)+x\cdot Du\,u_t+\frac{n-1}{2}uu_t,
\end{align*}
\begin{align*}
\mathbf{q}\triangleq\left(tu_t+x\cdot Du+\frac{n-1}{2}u\right)Du+\frac12(u_t^2-|Du|^2)x.
\end{align*}
将在 \S12.4 给出一个非线性波方程类似恒等式的重要应用。
\end{example}

\begin{example}[波方程的共形能量]\label{example:ch08-6-5-共形能量}
下面这个更为精致的例子说明，即使 Noether 定理不能直接应用，它有时仍可帮助识别有用的乘子。
映射
\begin{align}
(\overline{x},\overline{t})\triangleq\left(\frac{x}{|x|^2-t^2},\frac{t}{|x|^2-t^2}\right),\label{align:ch08-6-eq15}
\end{align}
其中 $|x|^2\neq t^2$，称为\textbf{双曲反演}。相关联的是双曲 Kelvin 变换 $\mathcal{K}u=\overline{u}$，定义为
\begin{align}
\overline{u}(x,t)\triangleq u(\overline{x},\overline{t})\big||\overline{x}|^2-\overline{t}^2\big|^{\frac{n-1}{2}}=u\!\left(\frac{x}{|x|^2-t^2},\frac{t}{|x|^2-t^2}\right)\frac1{\big||x|^2-t^2\big|^{\frac{n-1}{2}}}.\label{align:ch08-6-eq16}
\end{align}
第 12 章的一个习题表明：若 $\Box u=0$，则 $\Box\overline{u}=0$。

我们打算用双曲反演与双曲 Kelvin 变换帮助设计 $(x,t)$ 与 $u$ 中的变分。对前者，先把 $(x,t)$ 映到 $(\overline{x},\overline{t})$，再加 $\tau e_{n+1}$，最后再应用一次双曲反演。计算表明结果为映射
\begin{align}
\mathbf{x}(x,t,\tau)\triangleq\gamma(x,t+\tau(|x|^2-t^2)),\label{align:ch08-6-eq17}
\end{align}
其中
\begin{align}
\gamma\triangleq\frac{|x|^2-t^2}{|x|^2-(t+\tau(|x|^2-t^2))^2}.\label{align:ch08-6-eq18}
\end{align}
接着用类似程序构造 $u$ 的变分：对 $u$ 施加 Kelvin 变换算 $\mathcal{K}u=\overline{u}$，然后在 $\overline{u}$ 的自变量中加 $\tau e_{n+1}$，最后再施加 Kelvin 变换。另一计算显示所得函数变为
\begin{align}
w(x,t,\tau)\triangleq\gamma^{\frac{n-1}{2}}u(\mathbf{x}(x,t,\tau)).\label{align:ch08-6-eq19}
\end{align}
接着通过关于 $\tau$ 求导再令 $\tau=0$，计算对应于\eqref{align:ch08-6-eq17}--\eqref{align:ch08-6-eq19} 的乘子：
\begin{align}
\mathbf{v}=(2xt,|x|^2+t^2),\qquad m=(|x|^2+t^2)u_t+2tx\cdot Du+(n-1)tu.\label{align:ch08-6-eq20}
\end{align}
现我们并不断言作用泛函\eqref{align:ch08-6-eq14} 在这些定义域与函数变分下不变。反之，我们猜测：由于双曲 Kelvin 变换保持波方程的解，把波方程 $\Box u=0$ 乘以\eqref{align:ch08-6-eq20} 中的乘子 $m$ 也许有用。事实证明确实如此，经一个稍长的计算导出 Morawetz 恒等式
\begin{align}
c_t-\operatorname{div}\mathbf{r}=0,\label{align:ch08-6-eq21}
\end{align}
其中
\begin{align}
c\triangleq\frac12(|x|^2+t^2)(u_t^2+|Du|^2)+2tx\cdot Du\,u_t+(n-1)tu\,u_t-\frac{n-1}{2}u^2\label{align:ch08-6-eq22}
\end{align}
是所谓共形能量的密度，而
\begin{align}
\mathbf{r}\triangleq\big((|x|^2+t^2)u_t+2tx\cdot Du+(n-1)tu\big)Du+t(u_t^2-|Du|^2)x.\label{align:ch08-6-eq23}
\end{align}
Morawetz 恒等式之所以重要，是因为对 $n\neq 2$，共形能量密度 $c$ 可写成若干非负项之和加上一个 $x$ 变量中的散度：
\begin{align}
c=&\frac{(t+|x|)^2}{4}\left(u_t+u_r+\frac{n-1}{2|x|}u\right)^2+\frac{(t-|x|)^2}{4}\left(u_t-u_r-\frac{n-1}{2|x|}u\right)^2\notag\\
&+\frac{|x|^2+t^2}{2}\left(|Du|^2-u_r^2+\frac{(n-3)(n-1)}{4|x|^2}u^2\right)-\frac{n-1}{4}\operatorname{div}\!\left(\frac{|x|^2+t^2}{|x|^2}u^2x\right),\label{align:ch08-6-eq24}
\end{align}
其中 $u_r\triangleq\frac{Du\cdot x}{|x|}$。
\end{example}

\paragraph{应用：局部能量衰减。}现设 $O\subset\mathbb{R}^n$ 表示一个有界、光滑开集，关于原点星形：此条件的定义与几何意义见 \S9.4.2。定义外部区域
\begin{align*}
U\triangleq\mathbb{R}^n\setminus\overline{O}.
\end{align*}
另设 $u$ 是"障碍物" $O$ 外部下列波方程初边值问题的光滑解：
\begin{align*}
\begin{cases}
u_{tt}-\Delta u=0&\text{在 }U\times(0,\infty)\text{ 中},\\
u=0&\text{在 }\partial U\times\{t=0\}\text{ 上},\\
u=g,\ u_t=h&\text{在 }U\times\{t=0\}\text{ 上},
\end{cases}
\end{align*}
其中初值 $g,h$ 具有紧支集。

我们断言：若 $n\geqslant 3$ 且 $O\subset B(0,R)$，则存在常数 $C$ 使
\begin{align}
\int_{B(0,R)\setminus O}u_t^2+|Du|^2\,\mathrm{d}x\leqslant\frac{C}{t^2},\label{align:ch08-6-eq25}
\end{align}
对时间 $t\geqslant 2R$ 成立。于是任何有界区域内的能量随 $t\to\infty$ 衰减到零，尽管总能量守恒。

为证\eqref{align:ch08-6-eq25}，首先由守恒律\eqref{align:ch08-6-eq21} 观察到
\begin{align}
\frac{\mathrm{d}}{\mathrm{d}t}\int_U c\,\mathrm{d}x=\int_U\operatorname{div}\mathbf{r}\,\mathrm{d}x=\int_{\partial O}\mathbf{r}\cdot\nu\,\mathrm{d}S,\label{align:ch08-6-eq26}
\end{align}
其中 $\nu$ 表示 $\partial O$ 的内单位法向量。现 $u=u_t=0$ 在 $\partial O$ 上，故由\eqref{align:ch08-6-eq23} 可算得在 $\partial O$ 上
\begin{align*}
\mathbf{r}\cdot\nu=2t(x\cdot Du)(Du\cdot\nu)-t|Du|^2(x\cdot\nu).
\end{align*}
由于 $u=0$ 在 $\partial O$ 上，在那里 $Du=(Du\cdot\nu)\nu$。把此用于上式，推得
\begin{align*}
\mathbf{r}\cdot\nu=t|Du|^2(x\cdot\nu)\leqslant 0,
\end{align*}
因为 $O$ 关于原点星形且 $\nu$ 为内法向。

于是\eqref{align:ch08-6-eq26} 与关于 $c$ 的公式\eqref{align:ch08-6-eq24} 蕴涵：对每个时刻 $t>0$，
\begin{align*}
\int_{B(0,R)\setminus O}&\frac{(t+|x|)^2}{4}\left(u_t+u_r+\frac{n-1}{2|x|}u\right)^2+\frac{|x|^2+t^2}{2}\frac{(n-3)(n-1)u^2}{4|x|^2}\\
&+\frac{(t-|x|)^2}{4}\left(u_t-u_r-\frac{n-1}{2|x|}u\right)^2+\frac{|x|^2+t^2}{2}(|Du|^2-u_r^2)\,\mathrm{d}x\leqslant C.
\end{align*}
取 $t\geqslant 2R$ 并作简单估计，由此导出估计
\begin{align}
\int_{B(0,R)\setminus O}|Du|^2-u_r^2+\frac{(n-3)u^2}{|x|^2}\,\mathrm{d}x\leqslant\frac{C}{t^2}\label{align:ch08-6-eq27}
\end{align}
及
\begin{align}
\int_{B(0,R)\setminus O}u_t^2+u_r^2+\frac{n-1}{|x|}uu_r+\frac{(n-1)^2}{4|x|^2}u^2\,\mathrm{d}x\leqslant\frac{C}{t^2}.\label{align:ch08-6-eq28}
\end{align}
现
\begin{align*}
\frac{u}{|x|}u_r=\operatorname{div}\!\left(\frac{u^2}{2|x|^2}x\right)-\frac{n-2}{2}\frac{u^2}{|x|^2}.
\end{align*}
此恒等式表明：对 $n=3$，\eqref{align:ch08-6-eq28} 中最后两项的积分非负；对 $n>3$，可由\eqref{align:ch08-6-eq27} 估计它。能量衰减断言\eqref{align:ch08-6-eq25} 得证。

\end{document}
